% !TEX root = ../../../MathModel.tex
% 负责人：罗财能。
\subsubsection{具体分析}
\label{sec:q1:analysis}

第一问要求各架次沿$O01\to S_i\to O01$直接往返，只向一个服务区投送物资。
在不安排实体无人机和共享电池周转的前提下，决策包括单条航线的安全载荷、每架次使用的机型以及不可拆货箱的组合。
载重、舱容和返航余量共同限定可行组批；尤其在山区，航路中间山脊决定的爬升需求不能由起终点海拔替代。

本文先从原始DEM提取沿线最高地形，再建立载货去程、空载返程的能耗模型，求出各机型在各服务区的最大安全载荷。
随后，将同区同类且质量、体积相同的货箱压缩为计数状态，穷举全部可行组批模式，通过动态规划确定机型及逐架次货箱清单。
各区之间没有跨区装载或共享资源排程约束，因而可分区求解后合并。

原题要求说明架次、能耗和累计作业时间的权衡关系，但未指定唯一权重。
本文采用$\min_{\mathrm{lex}}(N,E,T)$：先最少架次，再在最少架次方案中选择最低能耗，最后比较累计作业时间。
减少架次可以减少起降和装卸组织次数；同时另求能耗优先与时间优先方案，定量检验这一选择的代价。
这里$T$为各架次作业时长之和，不是机队并行完工时间。
